% Commentary by GPT-6 Astra (OpenAI, ChatGPT Pro mode).

\begin{dcomment}{exposition}{p108:sur laquelle on peut se poser la question}{The total category and the telescope}
Grothendieck wants a path category that carries composition and at the same time has the intended homotopy type. The two objects in this opening comparison solve different parts of that problem. The total category remembers how the types of paths vary; the filtered colimit forgets the stages at which paths were represented. We write $BC=|NC|$, and call a functor a weak equivalence when its nerve is a weak homotopy equivalence. This is the meaning of an isomorphism in $\Hot$. It is stronger than a bijection on connected components and different from an isomorphism in $\mathrm{Cat\,hot}$, where morphisms are connected components of ordinary functor categories.

Here is a proof of the assertion announced on the following page. Let $F_0\to F_1\to\cdots$ be a strict sequence of small categories, let $E\to\mathbf N^o$ be its fibred Grothendieck construction, and let $L=\mathop{\rm colim}_nF_n$. The canonical functor $q:E\to L$ is a weak equivalence. For clarity, a morphism $(n,x)\to(m,y)$ in $E$, with $n\geq m$, is a morphism $x\to F_{m,n}y$ in $F_n$.

To prove the assertion, fix $z\in L$ and consider the comma category $q/z$. Any finite collection of its objects, arrows, and commuting triangles can be represented at one sufficiently large stage $r$: filtered colimits of categories are computed from finite diagrams and their eventual relations. We may also represent $z$ at stage $r$. Transporting the finitely many objects to that stage gives a diagram in $F_r/z_r$. Its projection to $q/z$ has a natural transformation to the original diagram, supplied by the cartesian arrows in $E$. It has a natural transformation to the constant diagram $(r,z_r,1)$, supplied by the arrows to $z_r$. Thus every finite simplicial diagram in $N(q/z)$ contracts after passing to a stage. The same argument with a chosen vertex proves connectedness and vanishing of all homotopy groups; the comma category is nonempty because $z$ has a representative. Quillen's Theorem A now gives the assertion \citep[\S1, Theorem A]{Quillen1973}. This argument does not require the transition functors to be faithful or injective on objects.

The second leg of the displayed roof is a separate matter. Take a nontrivial group $G$, let $T=BG$ be its one-object category, and take $C=T$ over itself. The preorder obtained by forgetting the distinction between parallel arrows is a point. A sequence constant on the identity arrow is therefore cofinal in that preorder, whereas the comparison is $*\to BG$. Its fundamental group is $1\to G$, so it is not a weak equivalence. This proves exactly the limitation noticed on manuscript page~109: cofinality in the preorder of objects does not account for the automorphisms of the indexing category. It does not predict failure for the particular category of interval types, which has additional structure.
\end{dcomment}

\begin{dcomment}{examples}{p112}{What the preorder forgets about a lift}
Grothendieck's proposed sufficient criterion fails for the category of types for a concrete reason: its preorder identifies distinct changes of parameter which cannot be connected inside a lifting category. The definition on manuscript page~99 retains the same objects in $\Ty_o$ and replaces every nonempty set of morphisms by a singleton. It does not identify different objects whose images happen to be isomorphic in that preorder.

Take
\[
 S=(0\longrightarrow1\longleftarrow2\longrightarrow3),
 \qquad T=(0\longrightarrow1).
\]
An endpoint-preserving monotone type map $S\to T$ changes its vertex value from $0$ to $1$ at exactly one of the two forward edges. There are precisely two such maps, with vertex lists
\[
                    (0,1,1,1),\qquad(0,0,0,1).
\]
Both $S$ and $T$ have only the identity endomorphism in $\Ty$. For $S$, a monotone endomorphism can advance by at most one vertex on each edge, since its alternating target has no nonidentity composite of two consecutive edges. Its three increments must sum to three because the endpoints are fixed, so each increment is one. The assertion for $T$ is immediate.

Let $p:\Ty\to\Ty_o$ and consider the unique base arrow $pS\to pT$. An incoming lift with target $T$ has source exactly $S$, and its morphisms are endomorphisms of $S$ commuting with the lifted arrow. Its lifting category is therefore the discrete category on the two displayed maps. The outgoing lifting category with source $S$ is likewise discrete on those maps, now because its morphisms are endomorphisms of $T$. Neither category is connected. The incoming and outgoing lift criteria consequently show that $p$ is neither $0$-smooth nor $0$-proper, and hence neither smooth nor proper. This verifies the source's assertion; failure of a sufficient criterion still does not imply failure of the geometric comparison itself.
\end{dcomment}

\begin{dcomment}{context}{p113:Je vais résumer les faits positifs}{The obstruction withdrawn in the footnote}
Grothendieck is trying to explain why an associative composition is easy on all finite paths but elusive on a chosen infinite refinement. The proposed explanation on this page is withdrawn in footnote~24, and the withdrawal is mathematically decisive. A based loop space does admit a strictly associative model. One elementary topological model consists of Moore loops: a loop is a pair $(\ell,\gamma)$, where $\ell\geq0$, $\gamma:[0,\ell]\to X$, and both endpoints are the basepoint. Concatenation adds lengths and joins the maps without rescaling their domains. Its associativity is an equality, and the loop of length zero is a strict unit. Reversal is an inverse up to homotopy, which is exactly the group law required in $\Hot$; a strict monoid model does not require strict inverses.

The simplicial statement is stronger. For a reduced simplicial set $K$, Kan's loop-group construction gives a simplicial group $GK$ and a weak equivalence $K\to\overline W GK$. Consequently $|GK|$ models the based loop space of $|K|$. We checked the precise reducedness hypothesis and the unit equivalence in \citet[Theorem~21]{Stevenson2012}, which gives a modern proof of Kan's theorem. The word ``semi-simplicial'' in this part of the manuscript includes the simplicial formalism with degeneracies; it must not be read as a claim about an arbitrary semisimplicial object without them.

These constructions settle the alleged obstruction, but do not by themselves identify either of Grothendieck's particular categories with the model. That identification is the work still to be done. For $\Sigma=\mathrm{Fl}(X)$, section~13 of the present fragment eventually compares the all-types category with $\Ch_\infty$ and restores the strict concatenation model. The footnote should therefore be read as a correction to an inference, not as an incidental qualification of a theorem. The unsuccessful search described in the surrounding prose is part of the discovery of the later proof.
\end{dcomment}

\begin{dcomment}{exposition}{p117:associant à tout $a \in X$}{What the ind-object represents}
The ind-object retains a uniform bound on the complexity of a family of paths. This is why it can encode localization before a single category has acquired a satisfactory geometric interpretation. With $[Y,D]_{\mathrm{cat}}=\pi_0\underline{\mathrm{Hom}}(Y,D)$, the formula is
\[
 H_T^Y(a,b)=\mathop{\rm colim}_n[Y,C_n(a,b)]_{\mathrm{cat}}
 \cong\mathrm{Hom}_{X(Y)[\Sigma_Y^{-1}]}(a_Y,b_Y).
\]
Here $\Sigma_Y$ consists of natural transformations whose components belong to $\Sigma$, and the right side uses the constant diagrams. A morphism on the right is represented by a finite zigzag of natural transformations. Its finite length is a bound shared by every object of $Y$; it therefore gives a map into one stage on the left. Refinement and the elementary relations of localization give the identifications in the colimit. This is the mechanism already discussed in the commentary on Chapter~V, manuscript pages~30--33, concerning the parameter category and ind-representation. The formula printed on page~117 omits the superscript $Y$ from $H_T^Y$; without it an ind-object is being equated to a set.

For finite $Y$, the natural map
\[
 \mathop{\rm colim}_n[Y,C_n]_{\mathrm{cat}}
       \longrightarrow [Y,\mathop{\rm colim}_n C_n]_{\mathrm{cat}}
\]
is bijective. Indeed a functor uses finitely many objects, arrows, and relations, all of which lift to one stage. A finite zigzag of natural transformations uses finitely many more data, so equality also occurs at one stage. Finite presentation, rather than any topological fibrancy, is doing the work. For infinite $Y$, a functor to the colimit can have an unbounded range of stages, and this argument stops.

Composition of localized zigzags is associative because juxtaposition of finite strings is associative. When all arrows of $X$ are inverted, reversal supplies the inverse string. Hence the family of sets $H_T^Y(a,b)$ is a groupoid with object set $\mathrm{Ob}(X)$. For a general $\Sigma$ it is only a category. Passing from this set-valued representation to mapping \emph{spaces} is an additional step: the notation does not identify the ordinary nerve of every functor category with a derived mapping object.
\end{dcomment}

\begin{dcomment}{development}{p120}{Contracting the category of choices}
Grothendieck seeks a coherent comparison between different cofinal sequences, so that changing a path parameter does not change the resulting operation in $\Hot$. His three-sequence diagram supplies this comparison, but the natural projections already imply a stronger statement about the space of choices.

\textbf{Contraction of the refinement category.}\ours\quad
Let $T$ be a category with a bifunctor $*$ and natural maps $t*t'\to t$ and $t*t'\to t'$. Let $\mathcal H$ be the nonempty category of sequences $\mathbf N^o\to T$ cofinal in the preorder of objects. Then $N\mathcal H$ is contractible. Consequently every functor $\mathcal H^o\to\mathcal A$ taking all arrows to isomorphisms has a unique system of comparison isomorphisms induced by arrows and their inverses; these comparisons depend only on their two endpoints and compose transitively.

\textit{Proof.} Fix $s\in\mathcal H$. Define $Q(t)=t*s$ pointwise, on both objects and morphisms of sequences. For every $a\in T$ there is an index $n$ and a map $t_n\to a$. Composing $t_n*s_n\to t_n\to a$ shows that $Q(t)$ is still cofinal. The two natural projections give natural transformations
\[
 \mathrm{id}_{\mathcal H}\ \longleftarrow\ Q\ \longrightarrow\ \mathrm{const}_s.
\]
The nerve of a natural transformation gives a homotopy of the associated maps on realizations: explicitly, a chain of arrows and a cut in that chain use the first functor before the cut, the transformation at the cut, and the second functor afterwards. Thus the identity of $B\mathcal H$ is homotopic to the constant map at $s$, proving contractibility.

A functor taking all arrows to isomorphisms factors through the localization of $\mathcal H$ at all its arrows, its fundamental groupoid. Contractibility makes this groupoid equivalent to a point: between any two of its objects there is exactly one morphism. Applying the functor to that morphism proves the uniqueness and transitivity assertions, including independence from the particular arrow of sequences used. The opposite variance changes the direction of the comparisons but not the argument.\hfill$\square$

Associativity of $*$ and a terminal unit are unnecessary for this proof once the two natural projections are given. In the interval application, concatenation has these projections by collapsing the other interval, and a cofinal sequence exists. The proposition isolates the formal part of the argument. It does \emph{not} prove that the actual functors $C_\infty(t')\to C_\infty(t)$ are weak equivalences; that premise requires the geometric comparison or a valid relative argument. In particular the conservativity assertion in footnote~26 cannot supply it in the stated generality. We checked the general nerve and cofinality mechanisms in Quillen's \S1 and the smooth-functor formulation of \citet[Theorem~1.2]{Maltsiniotis2005Asphericite}; we did not locate this explicit application to Grothendieck's category of sequences in those sources. The result is an elementary continuation of his construction, without a claim of priority.
\end{dcomment}

\begin{dcomment}{exposition}{p121}{Two different finiteness arguments}
The comparison on this page must keep finite categorical data separate from a finite homotopy type. The preceding block on ind-representation proves a filtered-colimit formula in $\mathrm{Cat\,hot}$ whenever the source category is finitely presented. In $\Hot$, the corresponding safe hypothesis is that the classifying space of the source has the homotopy type of a finite CW complex, or more generally is compact for the filtered homotopy colimits in question. A finite poset has a finite nerve and meets this hypothesis. But Grothendieck explicitly defines a finite category by finiteness of both its objects and its arrows in Chapter~V, manuscript page~45, and uses these categories to define $\Hot_0$ on pages~51--52 \citep{Grothendieck1990V}. With that convention, the assertion is false.

For a concrete distinction, let $Y$ be the one-object category of the group $\mathbf Z/2$. It is finite, but $BY\simeq\mathbf{RP}^{\infty}$. Consider the increasing skeleta $\mathbf{RP}^n$. The identity of $\mathbf{RP}^{\infty}$ cannot factor, up to homotopy, through one skeleton: such a factorization would make the identity on $H^k(\mathbf{RP}^{\infty};\mathbf F_2)$ factor through zero for every $k>n$, whereas these cohomology groups are nonzero. Thus
\[
 \mathop{\rm colim}_n[BY,\mathbf{RP}^n]
       \longrightarrow[BY,\mathbf{RP}^{\infty}]
\]
is not surjective. The same example can be expressed entirely in categories by using compatible face-poset models for the finite skeleta, or the categories of simplices of their simplicial models. It is the size of the nerve in all dimensions, not merely the number of arrows, that matters here.

For a finite simplicial source $K$, the valid filtered-colimit argument is direct. Replace the targets by a functorial Kan replacement built using finite horn data. Such a replacement commutes with filtered colimits. Maps from $K$ and simplicial homotopies between them involve finitely many simplices, so they appear at one stage. Finite CW sources follow by a finite simplicial model up to homotopy. The statement about filtered colimits of Thomason-fibrant categories is also valid, since the relevant horn domains are finite after applying the fixed subdivisions. It does not establish compactness for $B(\mathbf Z/2)$.

Accordingly, the literal last sentence of page~120 omits even the intended subscript zero; restoring that subscript would still leave the compactness error. A valid replacement restricts to homotopy-finite tests and thereby changes the stated test category. This repair does not alter the elementary finite-presentation calculation in $\mathrm{Cat\,hot}$.
\end{dcomment}

\begin{dcomment}{status}{p121}{The unpointed detection claim}
Grothendieck uses finite tests here to recognize a weak equivalence. In the unpointed homotopy category, the assertion that $\Hot\to\widehat{\Hot_0}$ is conservative is false. The editors already flag the claim in Chapter~V, manuscript page~52, and the commentary there gives the symmetric-group shift example; we do not repeat it. The general theorem is stronger: no set of spaces jointly detects all isomorphisms in the unpointed homotopy category \citep[Theorem~2.1]{ArlinChristensen2023}. Their Theorem~4.3 shows why retaining one more level of data helps: spheres do detect equivalences in the homotopy 2-category, whose morphism categories retain homotopy classes of homotopies.

This defect in the inference does not disprove the particular path comparisons. They can be proved with based or relative information, or, in the all-arrows situation of section~13, by comparing the explicit total path categories over $X\times X$. An audit must separate that geometric theorem from the invalid general detection principle used on this page. For based spaces, all basepoints and their homotopy groups give the usual Whitehead test. For free homotopy classes, the fundamental-group information lost to conjugacy must first be restored.
\end{dcomment}

\begin{dcomment}{examples}{p123}{A colimit need not be the same ind-object}
Grothendieck's warning in footnote~27 concerns the test objects allowed, not a defect of the colimit. It is already visible for sets. Put $D_n=\{0,\ldots,n\}$, viewed as discrete categories, and $D=\mathbf N$. For every finite discrete $K$ one has
\[
 \mathop{\rm colim}_n\mathrm{Hom}(K,D_n)=\mathrm{Hom}(K,D),
\]
because the image of a finite set is bounded. The map from the ind-object $(D_n)_n$ to the constant ind-object $D$ is nevertheless not an isomorphism. Applying maps from the constant object $\mathbf N$, its source consists of bounded-image functions $\mathbf N\to\mathbf N$, while the target also contains the identity. These computations are unchanged in $\Hot$: maps between discrete spaces are functions, and no homotopy identifies different functions.

Thus restriction to the finite-test category can identify two presheaves that differ on larger tests. The ind-object records how a family is bounded; the colimit records its pointwise values after the bound has been forgotten. This is the precise content behind ``pratiquement jamais essentiellement constant.'' It is not a universal prohibition: a sequence that is eventually constant is a constant ind-object, and more generally one can have an ind-isomorphism to a constant object without literal eventual constancy.
\end{dcomment}

\begin{dcomment}{exposition}{p126}{Padding at the two outer ends}
Grothendieck's construction now solves the binary-composition problem by choosing the locations where new constant segments are inserted. A common stage in each factor is harmless because the diagonal $\mathbf N\to\mathbf N\times\mathbf N$ is cofinal. The location of the padding is what makes the squares commute.

Write a path in the first system as $c_n$ and one in the second as $d_n$. Under $n\mapsto n+p$, the first becomes a string of $p$ constant blocks followed by $c_n$, while the second becomes $d_n$ followed by $p$ constant blocks. Concatenation gives
\[
 (1_a^{,p}*c_n)*(d_n*1_c^{,p})
       =1_a^{,p}*(c_n*d_n)*1_c^{,p},
\]
which is exactly the transition prescribed for the target system. The equality holds for the vertical transformations as well as for objects, because concatenation is a bifunctor. It therefore induces a functor on the filtered colimits. Products of categories commute with these filtered colimits, again because objects, arrows, and their defining equations are finite data.

A morphism $\tau_0\to\tau'_0$ gives a commutative square before passage to homotopy; once the refinement maps have been established to be weak equivalences and coherently identified, the square gives equality of the two binary operations in $\Hot$. To compare unrelated types, one can use $\tau_0*\tau'_0$ and its two collapse maps. The printed $\tau_0\times\tau'_0$ on this page cannot literally mean their categorical product as interval types: the product of two nontrivial intervals generally has a square rather than a linear underlying shape. Concatenation supplies the common refinement required by the argument.

The remaining associativity issue is real at the level of this proof. A third path has two joining positions, and the two binary constructions use different padding conventions at those positions. One must compare the resulting maps with coherent reparametrizations, or compare with a single strictly associative all-types model. Equality after restriction to finite tests cannot be substituted for that comparison. In the all-arrows case, the comparison established in section~13 is precisely the route back to the strict concatenation which motivated the discussion on page~111.
\end{dcomment}

\begin{dcomment}{examples}{p127}{A map invisible at every finite stage}
Grothendieck asks whether finite tests determine a morphism and whether the ind-object maps epimorphically onto its colimit. A small telescope gives a negative answer to both questions. We give the calculation because it exhibits the obstruction without relying on an unspecified phantom-map example.

Let $T$ be the mapping telescope of the degree-two maps
\[
 S^1\xrightarrow{2}S^1\xrightarrow{2}S^1\xrightarrow{2}\cdots.
\]
Its finite subtelescopes $T_n$ retract onto circles. Its fundamental group is $\mathbf Z[1/2]$, and its higher homotopy groups vanish: a map or nullhomotopy from a compact sphere lies in a finite subtelescope. Thus $T$ is a $K(\mathbf Z[1/2],1)$. After contracting the ray of vertices, a cellular model has one vertex, one loop $e_n$ for each $n\geq0$, and one two-cell with boundary $e_n-2e_{n+1}$. Cellular cohomology consequently gives
\[
 H^2(T;\mathbf Z)=\operatorname{coker}\left(
 \prod_{n\geq0}\mathbf Z\longrightarrow\prod_{n\geq0}\mathbf Z,
 \quad(b_n)\longmapsto(b_n-2b_{n+1})\right).
\]
This is the standard telescope computation of \citet[Examples~3F.7 and~3F.9]{Hatcher2002}.

To see nonvanishing directly, take $a_{2j}=1$ and $a_{2j+1}=0$. If $a_n=b_n-2b_{n+1}$, iteration would give
\[
 b_0=\sum_{j=0}^{N}2^ja_j+2^{N+1}b_{N+1}
 \qquad(N\geq0).
\]
In the ring of $2$-adic integers, this forces
$b_0=1+4+16+\cdots=-1/3$, impossible for an integer $b_0$. Hence $a$ determines a nonzero cohomology class. By the representation of integral degree-two cohomology, it determines an essential map
\[
 f:T\longrightarrow K(\mathbf Z,2)
\]
\citep[Theorem~4.57]{Hatcher2002}. Each restriction to $T_n$ is nullhomotopic, since $H^2(T_n;\mathbf Z)=0$. Every map from a finite CW complex into $T$ has image in a finite subtelescope up to homotopy, so all finite homotopy tests identify $f$ with the constant map.

Even if the test category includes nerves of all categories with finitely many objects and arrows, this example remains invisible. Each connected component of such a nerve has finitely generated fundamental group. Any homomorphism from that group to $\mathbf Z[1/2]$ has cyclic image, because finitely many denominators have a common power of two. The associated map to $T=K(\mathbf Z[1/2],1)$ therefore factors up to homotopy through $K(\mathbf Z,1)=S^1$, or through a point. Its composite with $f$ is nullhomotopic because $[S^1,K(\mathbf Z,2)]=0$.

Finally let $H=(T_n)_n$ in $\mathrm{Ind}\,\Hot$. The two maps $f,0:T\rightrightarrows K(\mathbf Z,2)$ become equal after precomposition with $H\to T$, since
\[
 \mathrm{Hom}_{\mathrm{Ind}\,\Hot}(H,K(\mathbf Z,2))
      =\mathop{\rm lim}_n[T_n,K(\mathbf Z,2)]=\{0\}.
\]
They are distinct before precomposition. Hence $H\to T$ is not an epimorphism. This example refutes the proposed general inference; it does not contradict associativity of the particular path composition, which requires its own geometric proof.
\end{dcomment}

\begin{dcomment}{exposition}{p128}{Edges and the fibres of the path object}
Grothendieck wants the categorical construction to work before fibrant replacement. The distinction needed here concerns the endpoint fibres. For a Kan complex $K$, the map $K^{\Delta^1}\to K\times K$ is a Kan fibration, and its fibre over two vertices models their space of paths. For an arbitrary simplicial set, the same strict fibre need not compute a homotopy fibre. This does not mean that the total space $K^{\Delta^1}$ has the wrong homotopy type: evaluation at an endpoint has the constant-path section and the simplicial interval contraction gives a homotopy from their composite to the identity.

A finite example makes the endpoint problem explicit. Let $X$ be the category with objects $0,1,2$ and nonidentity arrows $0\to1\leftarrow2$. Its classifying space is contractible, so the topological space of paths from $0$ to $2$ is nonempty and contractible. But there is no morphism $0\to2$. The fibre of the functor category $\underline{\mathrm{Hom}}([1],X)\to X\times X$ over $(0,2)$ is therefore empty. The two-edge zigzag already supplies the missing path, and the refinement formalism is designed to supply all the further homotopies between such representatives.

There is a related distinction in the discussion of composition. Two composable edges of a simplicial set specify the spine $\Delta^1\cup_{\Delta^0}\Delta^1$, not an entire $\Delta^2$. For a Kan complex a filler produces a third edge, with higher fillers controlling the dependence on choices. In an ordinary category the interval $[2]$ already includes the composite arrow; this explains why the categorical notation for two direct segments can conceal the simplicial filling step. A simplicial loop group, a Moore-path model, or a coherent localization retains the data needed to compare repeated compositions. None requires an algorithm computing all homotopy groups of a finite simplicial set.

The last conjecture on this page is correct for the all-arrows path object: a weak equivalence of categories induces weak equivalences on the corresponding derived path spaces, with their endpoints carried along. The justification for identifying Grothendieck's specific categories with those spaces belongs to the geometric argument in sections~12--13, not to a calculation with one-edge paths.
\end{dcomment}

\begin{dcomment}{development}{p121}{A comparison that fails under the general hypotheses}
The preceding distinction can be tested inside the exact fibred setting of section~10, rather than only in the abstract homotopy category. Grothendieck assumes there that parallel reindexing functors are homotopic, that the preorder of the base is cofiltered and has a countable coinitial subset, and that a strict cleavage has been chosen. These assumptions produce the same ind-object, but they do not force every map between the resulting colimits to be a weak equivalence.

\textbf{A fibred symmetric-group example.}\ours\quad
There is a strict fibred category $C\to T$ satisfying these assumptions, with $T$ also admitting an associative product with terminal unit and natural projections, for which a morphism of cofinal sequences induces an isomorphism in $\mathrm{Ind}\,\mathrm{Cat\,hot}$ and a noninvertible map in $\Hot$.

\textit{Proof.} Let $\mathrm{FI}$ be the skeleton of finite sets and injections, with objects $[n]=\{0,\ldots,n-1\}$, and set $T=\mathrm{FI}^o$. Write $\mathbf B G$ for the one-object category of a group $G$. Define a strict functor
\[
 F:T^o=\mathrm{FI}\longrightarrow\mathrm{Cat},\qquad
 F([n])=\mathbf B\mathfrak S_n.
\]
For an injection $j:[n]\hookrightarrow[m]$, send a permutation $\sigma$ to the permutation $j\sigma j^{-1}$ on $j([n])$, extended by the identity on its complement. Extension along a composite of injections agrees exactly with successive extension, so $F$ is a functor. Its Grothendieck construction is strictly fibred over $T$.

If $j,k:[n]\hookrightarrow[m]$ are parallel injections, choose $h\in\mathfrak S_m$ with $h j=k$, by extending the resulting bijection of their images to the complements. Then
$F(k)(\sigma)=hF(j)(\sigma)h^{-1}$ for every $\sigma$. Conjugation by $h$ is a natural isomorphism between the associated one-object functors. Thus the homotopy hypothesis is satisfied. The preorder $T_o$ is $\mathbf N^o$. Disjoint union provides the required associative operation and the empty set is its terminal unit in $T$; the two inclusions in $\mathrm{FI}$ become its natural projections.

Take $t'_n=[n]$ and $t_n=[n+1]$, with the opposite of the usual initial-segment inclusions as transitions. Both sequences are cofinal. The injections
$\delta_n:[n]\hookrightarrow[n+1]$, $i\mapsto i+1$, commute with those inclusions and therefore define a morphism $t\to t'$ in $\mathcal H$. The two filtered colimits of the fibres identify with $\mathbf B G$, where $G$ is the group of finite-support permutations of $\mathbf N$. The induced functor is the homomorphism
\[
 s:G\longrightarrow G,\qquad
 (s\sigma)(0)=0,\quad(s\sigma)(i+1)=\sigma(i)+1.
\]
It is not surjective, since the transposition $(0\ 1)$ is absent from its image. Consequently $Bs$ does not induce an isomorphism on fundamental groups and is not a weak equivalence.

At each finite stage, however, $\delta_n$ and the standard inclusion are parallel injections, so their induced functors are naturally isomorphic. In $\mathrm{Ind}\,\mathrm{Cat\,hot}$ the comparison is therefore the ordinary cofinal-shift isomorphism. This proves all assertions.\hfill$\square$

The earlier Chapter~V counterexample supplies the shift mechanism; the additional content here is its realization by a strict fibred category satisfying the hypotheses actually used on pages~117--123. We checked those hypotheses against the editors' complete Chapter~VII, manuscript pages~99--101, and the detection theorem of \citet{ArlinChristensen2023}; we did not find this indexed realization in those sources. No priority is claimed. The example does not satisfy the stronger full-faithfulness features of the geometric padding maps. It leaves the later all-arrows path theorem intact, while disproving the general inference that an ind-isomorphism of homotopy classes must remain an isomorphism after these ordinary colimits are formed.
\end{dcomment}




\begin{dcomment}{exposition}{p129:Rapide retour sur les}{What an external Hom must represent}
Grothendieck now asks what his paths would have to accomplish in order to serve a theory of derived diagrams. A path from $a$ to $b$ is no longer only a point of a topological mapping space. It must also describe a morphism between the constant diagrams $a_Y$ and $b_Y$, after objectwise localization of $Y$-diagrams. This is the reason for testing the path construction against every parameter category $Y$.

Write
\[
 \mathbb D(Y)=\operatorname{Fun}(Y,\mathcal M)[\Sigma_Y^{-1}],
 \qquad
 F_{a,b}(Y)=\operatorname{Hom}_{\mathbb D(Y)}(a_Y,b_Y).
\]
The conventions on variance can be reversed together without changing this discussion. The first displayed condition says precisely that $F_{a,b}$ factors through the localization of categories at nerve weak equivalences. It is a necessary invariance property for the proposed representation by a space. It is not by itself a definition of a derivator: existence and pointwise computation of homotopy Kan extensions are additional requirements. For combinatorial model categories these requirements hold by \citet[Proposition~1.30]{Groth2013}. For their locally constant diagrams, the invariance asserted here is a theorem of \citet[Theorem~1.3]{Cisinski2009LocallyConstant}; constant diagrams are particular locally constant diagrams, and the subcategory in that theorem is full.

The contemporary form of the intended answer is especially transparent when a relative category presents an $\infty$-category $\mathcal C$ and its localized strict diagram categories present $\operatorname{Fun}(NY,\mathcal C)$. With $K=\operatorname{Map}_{\mathcal C}(a,b)$ one then has
\[
 F_{a,b}(Y)\cong
 \pi_0\operatorname{Map}_{\mathcal S}(|NY|,K).
\]
Indeed, a coherent transformation between the two constant diagrams is a coherent $Y$-indexed family of points of $K$; this is the limit of the constant $K$-diagram, hence the mapping space from the groupoid completion $|NY|$ to $K$. Taking components gives the formula. The diagram-presentation hypothesis is substantive: one cannot replace it by the existence of an arbitrary admissible set $\Sigma$. The comparison between strict model-category diagrams and coherent diagrams is established, under the relevant model-category hypotheses, in \citet[Proposition~4.2.4.4]{Lurie2009HTT}.

Here the source means finite sets of objects and arrows by a finite category \citep[ms.~45]{Grothendieck1990V}; such a category need not have a finite nerve. Consequently the earlier compactness argument does not establish the comparison for every finite category. The conditional assertion under the additional locality of $H$ is not disproved merely by that observation.

This also clarifies the finite-stage issue. The ind-object $H(a,b)$ records a bound on the length of a zigzag that works simultaneously over the entire parameter. Passing to $\Ch_\infty$ can forget this uniformity. It does not follow that modern derived mapping spaces must themselves be ind-objects: in the good diagram presentations just described, one space $K$ represents the functor for every small $Y$. What fails for a general pair $(\mathcal M,\Sigma)$ may be the identification of its localized strict diagrams with coherent diagrams, rather than representability of the latter. The example below separates these two issues with a six-object category.
\end{dcomment}

\begin{dcomment}{context}{p131}{The original category and its localized value}
Grothendieck wants the same external Hom to explain both homology and cohomology with constant coefficients. This is an adjunction calculation once the ambient categories have been fixed. Put $p:Y\to e$. If the relevant adjoints exist, then
\[
\begin{aligned}
 \operatorname{Hom}_{\mathbb D(Y)}(p^*a,p^*b)
 &\cong\operatorname{Hom}_{\mathbb D(e)}(p_!p^*a,b)\\
 &\cong\operatorname{Hom}_{\mathbb D(e)}(a,p_*p^*b).
\end{aligned}
\]
Here $\mathbb D(e)=\mathcal M[\Sigma^{-1}]$, not the original $\mathcal M$ unless the notation has been changed. Likewise the $a_X,b_X$ of the displayed formula mean $a_Y,b_Y$. These small notational adjustments matter because an adjoint to restriction of localized diagrams is a derived operation; an adjoint before localization need not compute it.

The historical motivation is unusually explicit. In the retrospective in \emph{Pursuing Stacks}, \citet[\S69, pp.~1--2]{Grothendieck1983} relates the inadequacy of an isolated derived category to the trace questions arising from SGA~5, and proposes retaining derived categories of all diagram shapes. The present pages carry that programme into a concrete question: can the dependence on the parameter shape itself be represented homotopically? The adjunction formulas explain why a positive answer belongs simultaneously to the homological and the cohomological sides of the theory. The subsequent calculations concern the geometric adequacy of the proposed representative, rather than another formal manipulation of these adjunctions.
\end{dcomment}

\begin{dcomment}{development}{p134}{A fixed circle already sees every degree}
Grothendieck's intuition that localization of strict diagrams can depend on the chosen category presenting a homotopy type is correct. The proposed explanation by a bound on the degree of a four-arrow circle is not. A natural transformation between long paths contains an arbitrarily wide commuting ladder. Its contribution to a two-dimensional homotopy class is not bounded by the number of arrows in the parameter category. We give a finite example that both disproves the suggested bound and proves the desired failure of invariance.

\textbf{A suspension calculation.}\ours\ Let $Y$ be the four-object poset with minimal objects $r_0,r_1$, maximal objects $q_0,q_1$, and all four relations $r_i<q_j$. Let $X$ be obtained by adjoining two incomparable objects $s,n$, each greater than every object of $Y$. Take all arrows of $X$ as weak equivalences. Then $|NY|$ is a circle and $|NX|$ is its suspension, hence a $2$-sphere. For the constant diagram $s_Y$ the group
\[
 G_Y=\operatorname{Aut}_{\Pi\operatorname{Fun}(Y,X)}(s_Y)
\]
contains an element of infinite order. Moreover, for one fixed four-arrow $Y$, functors $Y\to\Ch_\infty(X;s,s)$ realize every integer in $\pi_1\Omega|NX|$.

\emph{Proof.} Denote the inclusion by $j:Y\hookrightarrow X$, and choose a vertex $y_0\in Y$. Since all objects of $Y$ lie below both suspension vertices, there is a loop of natural transformations
\[
 s_Y\ \longleftarrow\ j\ \longrightarrow\ n_Y
       \ \longleftarrow\ (y_0)_Y\ \longrightarrow\ s_Y.       \tag{B1}
\]
This is a loop in the nerve of the functor category, based at $s_Y$. Under evaluation it gives a continuous family of loops in $|NX|$: at $y$ the loop travels from $s$ to $n$ along the meridian through $y$, and returns along the fixed meridian through $y_0$.

To check its class, let $\gamma$ be an oriented fundamental $1$-cycle of $|NY|$. The two varying portions of the family sweep out the two cones on $\gamma$. The fixed return meridian contributes no $2$-chain. Thus the image of the fundamental class of $|NY|\times S^1$ is, with a choice of orientation,
\[
                    n*\gamma-s*\gamma,
\]
which is the fundamental cycle of the suspension $|NX|$. In particular its degree is $1$. Equivalently, collapse the fixed meridian in the suspension: the adjoint of the family becomes the usual suspension map of the identity of $|NY|$. The collapse of that contractible arc induces an isomorphism on $H_2$, so this gives the same calculation. Consequently (B1) is nontrivial and its $k$-fold power has degree $k$. A null-homotopy in the functor category would give a null-homotopy after evaluation and degree zero. All these powers are therefore distinct in $G_Y$.

The same construction takes place in one finite path stage. For
$\rho=\Delta_-\Delta_+\Delta_-\Delta_+$ define
\[
 F(y)=\bigl(s\longleftarrow y\longrightarrow n
                    \longleftarrow y_0\longrightarrow s\bigr).
\]
For an arrow $y\to y'$ use that arrow in the varying position and identities in the other four positions. All squares commute because $X$ is a poset. This defines a functor $F:Y\to\Ch_\rho(X;s,s)$. Concatenating $k$ copies pointwise gives a functor into $\Ch_{\rho^k}(X;s,s)$ of degree $k$. Insert identity arrows at the two ends to fit this alternating word into a sufficiently long symmetric stage $t_N$; this changes only parametrization, not the degree. Reversal gives negative $k$, and the constant family gives zero. This proves the assertion for $\Ch_\infty$, without assuming that its concatenation has already been made into a globally defined bifunctor.\hfill$\square$

The example also supplies a genuine change-of-parameter failure. Write $B\mathbb Z$ for the one-object category with automorphism group $\mathbb Z$. There is a functor $u:Y\to B\mathbb Z$ sending one of the four edges to $1\in\mathbb Z$ and the other three to $0$. Choose the sign so that the induced map on the circle has degree $1$. Its nerve is a weak equivalence: both classifying spaces are $K(\mathbb Z,1)$ and the induced map on fundamental groups is an isomorphism. Every functor $B\mathbb Z\to X$ is constant, since a poset has only identity automorphisms; every natural transformation between such functors is exactly one arrow of $X$. Therefore
\[
 \operatorname{Fun}(B\mathbb Z,X)\cong X,
 \qquad
 \operatorname{Aut}_{\Pi\operatorname{Fun}(B\mathbb Z,X)}(s)=1.
\]
Restriction along $u$ maps this trivial group to $G_Y$, which contains all the powers of (B1). It is not surjective. By the localization interpretation of $H(s,s)$ used on ms.~129, this is precisely a failure of the required invariance of $H(s,s)$.

We checked the relevant diagram-localization and fibrancy results in \citet{Cisinski2009LocallyConstant}, \citet{Thomason1980}, and \citet{Raptis2010}. We did not find this finite suspension calculation presented there. It is included as a direct completion of Grothendieck's proposed example, without a claim of priority. The failure proved here concerns localized \emph{strict} diagrams in this relative category. It is consistent with the homotopy invariance of coherent diagrams explained above.
\end{dcomment}

\begin{dcomment}{examples}{p135:Mais la flèche $(*)$ n'a aucune envie}{The absolute arrow category and its endpoint fibres}
Grothendieck is trying to isolate the defect that forced him to lengthen paths. The defect is real, but the specific absolute comparison on this page does not exhibit it. For every category $C$, put $A=\operatorname{Fun}([1],C)$, let $c:C\to A$ send an object to its identity arrow, and let $e_0:A\to C$ evaluate at the source. Then $e_0c=\mathrm{id}_C$. For an arrow $f:x\to y$, the square with vertical arrows $\mathrm{id}_x$ and $f$ defines a natural transformation
\[
                    ce_0\Longrightarrow\mathrm{id}_A.
\]
After taking nerves, this is a homotopy from $Bce_0$ to the identity. Hence $Be_0$ is a homotopy equivalence. Evaluation at $0$ also gives a homotopy equivalence
$\operatorname{Map}(|N[1]|,BC)\to BC$, since the interval contracts to $0$. The canonical comparison
\[
 B\operatorname{Fun}([1],C)
       \longrightarrow\operatorname{Map}(|N[1]|,BC)
\]
commutes with these evaluations and is consequently a weak equivalence. No fibrancy assumption on $C$ is involved.

What cannot be inferred is the corresponding assertion on \emph{strict} endpoint fibres. The fibre over $(a,b)$ of $A\to C\times C$ is the discrete category $\operatorname{Hom}_C(a,b)$: a square with identity arrows at its two ends can only relate an arrow to itself. The homotopy fibre of the topological path fibration is instead the full path space from $a$ to $b$. The fork example discussed at ms.~128 already shows that these can differ even on components. Thus the motivation for alternating paths survives, while the sentence claiming failure of the absolute $Y=[1]$ comparison must be corrected.

There is also a direct example of absolute failure for another shape. Let $P$ be any poset with $BP\simeq S^1$, and take $Y=B\mathbb Z$. As above, $\operatorname{Fun}(B\mathbb Z,P)\cong P$. Its classifying space is connected, whereas $\operatorname{Map}(S^1,S^1)$ has one component for every integer degree. The comparison misses every nonconstant degree. This pinpoints the missing derived dependence on the source, without confusing the total arrow category with its fibres.
\end{dcomment}

\begin{dcomment}{development}{p136:On voit donc que les catégories}{A fibrant target does not derive every source}
Grothendieck's next attempt is to impose a Kan condition on the target. His footnote identifies one missing assertion about products of cofibrations, but the displayed proof also replaces arbitrary maps in the localized category by ordinary functors from an arbitrary source. Thomason fibrancy alone does not permit that replacement. The following example makes the issue visible already on components.

\textbf{A fibrant-poset obstruction.}\ours\ There are a Thomason-fibrant category $P$ and a finite category $Y$ for which
\[
 B\operatorname{Fun}(Y,P)\longrightarrow
                 \operatorname{Map}(BY,BP)
\]
is not a weak equivalence. Here ``finite'' means finitely many objects and arrows, not a finite nerve.

\emph{Proof.} Let $G=\mathbb Z/2$ and $Y=BG$. By \citet[Theorem~2.6]{Raptis2010}, Thomason's model structure restricts to posets, with the same weak equivalences and fibrations. Start with a poset representing $BG$, for example the twofold subdivision construction $c\operatorname{Sd}^2N(BG)$, and take its fibrant replacement within posets. We obtain a poset $P$ that is fibrant as a category and satisfies $BP\simeq BG$. Every functor $BG\to P$ sends the single object to some $p$ and every element of $G$ to $\mathrm{id}_p$. Natural transformations are exactly the relations $p\leq p'$. Hence
\[
                        \operatorname{Fun}(BG,P)\cong P.
\]
The source of the comparison is therefore connected. Its image consists of constant maps $BG\to BP$. But the target has at least two components: a constant map induces the zero homomorphism on $\pi_1$, whereas a homotopy equivalence $BG\to BP$ induces an isomorphism. These maps cannot be homotopic, even without specified basepoints, since conjugation cannot turn a trivial homomorphism into an isomorphism. The comparison fails on $\pi_0$.\hfill$\square$

This also supplies an immediate test of the first isomorphism asserted in the proof on ms.~136: take its $Z=e$. The claimed bijection would identify the connected set $\pi_0\operatorname{Fun}(Y,P)$ with $[BY,BP]$, which has the two distinct classes just exhibited. One cannot repair that step merely by proving that the functor category is fibrant. In Thomason's model structure the source has to be handled as well; in particular every cofibrant category is a poset \citep[Proposition~5.7]{Thomason1980}, so this $BG$ is not cofibrant.

The literal product claim in footnote~29 would also be too strong for arbitrary cofibrations: multiplying the cofibration $\varnothing\to e$ by $Y$ would make every category cofibrant. The next development also rules out preservation of all trivial cofibrations under product with arbitrary $Y$: the lifting adjunction would otherwise force every such functor category to be fibrant.

A safe model for the desired derived mapping space is the simplicial function complex
\[
              \operatorname{Map}_{\mathrm{sSet}}(NY,RNP),
\]
where $RNP$ is a Kan replacement. Every simplicial set is cofibrant, and this construction derives the target in the cartesian model category of simplicial sets. The ordinary nerve identity
$N\operatorname{Fun}(Y,P)=(NP)^{NY}$ explains precisely what the uncorrected formula omits. We checked the model structure and the poset replacement in the cited papers; we did not find this particular $B(\mathbb Z/2)$ application there. The result is recorded as an elementary application, not a priority claim.
\end{dcomment}

\begin{dcomment}{development}{p136}{Fibrancy is not preserved by strict diagram categories}
The failure of the preceding mapping-space formula does not logically decide the separate assertion that $\operatorname{Fun}(Y,C)$ is fibrant whenever $C$ is. That assertion also fails, but it needs a different example. An equivariant construction makes the obstruction precise: ordinary horn filling need not respect an action on the diagram.

\textbf{A finite parameter with a nonfibrant functor category.}\ours\quad
There are a countable Thomason-fibrant category $E$ and a finite category $Y$ such that $\operatorname{Fun}(Y,E)$ is not Thomason-fibrant.

\emph{Proof.} Let $D$ be the four-object crown used above, so $BD\simeq S^1$. It is not Thomason-fibrant. Indeed, $\operatorname{Ex}^2ND$ has finitely many $1$-simplices: they are the functors from the finite poset $c\operatorname{Sd}^2\Delta^1$ to $D$. In a Kan complex every element of a based fundamental group is represented by one loop edge, since horn fillers compose and invert edge paths. If $\operatorname{Ex}^2ND$ were Kan, its fundamental group would therefore be finite, contrary to its being isomorphic to $\pi_1(BD)=\mathbb Z$.

Put $G=\mathbb Z/2$ and give $D$ the trivial action. We first construct a $G$-poset $Q$ whose underlying category is fibrant and whose fixed subposet is exactly $D$. Start with $Q_0=D$. At step $r$, for every generating Thomason horn inclusion
\[
 A=c\operatorname{Sd}^2\Lambda_i^n
       \ \longrightarrow\ B=c\operatorname{Sd}^2\Delta^n
\]
and every ordinary functor $F:A\to Q_r$, attach the free $G$-cell $G\times B$ along $G\times A$, using the equivariant attaching map $(g,a)\mapsto gF(a)$. Here $G$ in a product denotes the discrete set of its two elements. Take all these pushouts simultaneously and put $Q=\varinjlim_rQ_r$.

These horn inclusions are Dwyer inclusions of finite posets. The pushout description in \citet[Lemmas~4.7--4.8]{MayZakharevichStephan2017} shows that the old poset stays full and that the new objects are exactly the disjoint sets $G\times(B\setminus A)$. All new objects have free $G$-action. Thus $Q^G=D$, with its original order. Every horn map into $Q$ factors through some finite stage, because its domain is finite, and it obtains a filler at the next stage. Consequently $Q$ is Thomason-fibrant. There are only countably many finite horn shapes and, at each countable stage, countably many maps from them. Hence $Q$ is countable.

Define the transport category $E=Q\rtimes G$ as follows. Its objects are those of $Q$, and an arrow $x\to y$ labelled $g\in G$ exists exactly when $gx\leq y$; composition multiplies labels. There is a projection $E\to BG$. We claim that $E$ is Thomason-fibrant. Given a horn diagram $F:A\to E$, write $g_{ab}$ for the label of an arrow $a\to b$. Since $BA$ is contractible, its fundamental groupoid is equivalent to a point. The projected functor $A\to BG$ is therefore naturally isomorphic to a constant functor: choose labels $h_a$ with
\[
                         h_b g_{ab}=h_a.
\]
Replacing $F(a)$ by $h_aF(a)$ turns every arrow label into $1$ and gives a functor $A\to Q$. Extend it to $B\to Q$ by fibrancy of $Q$. Extend the assignments $h_a$ arbitrarily to the remaining objects of $B$, and undo the replacement: at $b$ use $h_b^{-1}F'(b)$, and label $b\to c$ by $h_c^{-1}h_b$. The inequalities defining $F'$ make this a functor $B\to E$, and it restricts exactly to the original horn diagram. This proves the claim.

Finally set $Y=BG$. A functor $BG\to E$ induces a homomorphism $\phi:G\to G$, either trivial or the identity. The trivial case is the identity action at an arbitrary object of $E$. In the identity case, an endomorphism labelled by the nontrivial $g$ requires $gx\leq x$; applying $g$ gives the reverse inequality, so $gx=x$. Its objects are therefore precisely $D=Q^G$. A natural transformation labelled by $h$ must satisfy $h\phi(g)=\psi(g)h$. Since $G$ is abelian, this forces $\phi=\psi$. For the identity case the remaining condition is just $x\leq y$, and every $h\in G$ is allowed. We obtain an actual decomposition
\[
             \operatorname{Fun}(BG,E)
                      \cong E\amalg(D\times BG).
\]
The category $D$ is a retract of this coproduct: project the second summand onto $D$, send the first constantly to a chosen object of $D$, and use identity labels for the section. A retract of a fibrant category is fibrant. Since $D$ is not fibrant, neither is $\operatorname{Fun}(BG,E)$.\hfill$\square$

Thus both claims on ms.~136 fail: arbitrary strict diagram categories need not preserve Thomason fibrancy, and even when a particular functor category is fibrant, it need not represent the derived mapping space. The equivariant cell and fixed-point ingredients above are standard and are cited at their use. We did not locate this application to strict diagram categories in the primary sources checked, including the published equivariant-poset paper; the argument is supplied here without a claim of priority.
\end{dcomment}

\begin{dcomment}{exposition}{p139:Voici pour le moment les propriétés essentielles}{The two maps in a mapping-path factorization}
Grothendieck turns from representing an external Hom to replacing a map by its homotopy fibres. This change of purpose makes the endpoints decisive. Put $E=\Ch_\infty(X)$, with source $\sigma$, target $\beta$, and constant-path section $i$. Then
\[
 C(f)=X'\times_{X,\sigma}E,
 \quad r:C(f)\longrightarrow X',
 \quad \bar f=\beta\operatorname{pr}_E,
 \quad \alpha(x')=(x',i(fx')).
\]
The strict identities are $r\alpha=\mathrm{id}$ and $\bar f\alpha=f$. Geometrically, $r$ forgets a path whose initial point is prescribed; its fibre should be contractible. The other map remembers the final point; its fibre should be the homotopy fibre of $f$. Thus the weak equivalence is $\alpha$, whereas $\bar f$ is a weak equivalence exactly when $f$ is. The sentence on ms.~139 that attaches the unconditional weak-equivalence assertion to $\bar f$ is contradicted by item~$3^o$ immediately below it.

Once $r$ is known to be a weak equivalence, the rest of this assertion uses only two-out-of-three. Since $r\alpha=\mathrm{id}$, the section $\alpha$ is a weak equivalence. Since $\bar f\alpha=f$, the maps $f$ and $\bar f$ are weak equivalences simultaneously. Neither inference says that $\alpha$ is a Thomason cofibration, or that $\bar f$ is a Thomason fibration. Smoothness and properness concern the homotopy exactness of base change, while Thomason's chosen fibrations are tested by $\operatorname{Ex}^2N$. Their definitions must not be interchanged.

The cofibration expectation actually fails for the unaltered constant-path inclusion. Take $f=\mathrm{id}_X$ with $X=[1]$. A nonconstant alternating path from $0$ to $1$, padded by identity arrows, admits a natural transformation to the constant path at $1$: at a vertex labelled $0$ use $0\to1$, and at a vertex labelled $1$ use its identity. Hence the image of $\alpha:X\to E$ is not a sieve. Every Thomason cofibration is a pseudo-Dwyer map and therefore a sieve \citep[\S2, following Proposition~2.4]{Raptis2010}. Thus $\alpha$ need not be a Thomason cofibration. This does not damage the intended homotopy-fibre replacement; it limits the proposed identification with a particular model-category factorization.

The expectation that $\bar f$ is always a Thomason fibration also fails directly. Take $X=e$. Since $\Ch_\infty(e)=e$ strictly, the construction gives $C(f)=X'$ and $\bar f=f:X'\to e$. Choosing $X'=D$, the nonfibrant crown of the preceding development, gives a counterexample. The homotopical lifting properties proved later do not impose Thomason fibrancy on an arbitrary category mapping to $e$.

The properties in items~$1^o$--$3^o$ are requirements at this stage of the manuscript. Their proofs begin only after the concrete symmetric model is fixed on ms.~145. We will use them in that order, rather than reading the list as an already established theorem for an unspecified projective system.
\end{dcomment}

\begin{dcomment}{exposition}{p140:Donc on trouve, à isomorphisme canonique près}{Two different reversals}
Grothendieck needs two symmetries because two different operations are being compared: exchanging the endpoints of a path, and replacing the ambient category by its opposite. For a signed word $\tau=(\varepsilon_1,\ldots,\varepsilon_r)$, with $\varepsilon_i\in\{+,-\}$, write
\[
 \tau^o=(-\varepsilon_r,\ldots,-\varepsilon_1),
 \qquad
 \tau^-=(-\varepsilon_1,\ldots,-\varepsilon_r).
\]
The first operation reverses the ordered interval. It is an anti-automorphism of concatenation and interchanges source and target without changing $X$. The second reverses every categorical arrow while preserving the ordering of the vertices. It is an automorphism of concatenation and is the operation accompanying $X\mapsto X^{\mathrm{op}}$.

At a fixed stage the corresponding identities are obtained by reindexing diagrams and natural transformations:
\[
 \Ch_\tau(X)\cong\Ch_{\tau^o}(X),
 \qquad
 \Ch_\tau(X^{\mathrm{op}})
          \cong\Ch_{\tau^-}(X)^{\mathrm{op}}.
\]
The first identity is understood over the transposition of $X\times X$; the second is over $(X\times X)^{\mathrm{op}}$. Passing through the specified transition maps gives the identities for projective systems. For $\tau_0=\Delta_+\Delta_-$ one has $\tau_0^o=\tau_0$ but $\tau_0^-=\Delta_-\Delta_+$. This explains why path reversal becomes a strict involution in the symmetric model, while categorical duality initially changes its parity. The translation introduced on ms.~147 supplies the remaining comparison. The commentary at manuscript p.~9 already distinguishes these operations; here their role is to transfer the endpoint and base-change assertions.
\end{dcomment}

\begin{dcomment}{later}{p142:Ou qu'on ait tout au moins des isomorphismes canoniques}{Transport and the comparison with ordinary path spaces}
Grothendieck wants the replacement to remember how its fibres vary, and his marginal request to compute $\pi_0$ is the right concrete test. For the ordinary mapping-path construction, a point in the homotopy fibre at $a$ is a pair $(x',\gamma)$ with $\gamma:f(x')\leadsto a$. Transport along $c:a\leadsto b$ sends it to
\[
                       (x',\gamma*c).
\]
Homotopic paths give homotopic maps of fibres; concatenation and reversal give composition and inverse in the homotopy category. This produces the functor from the fundamental groupoid requested in item~$4^o$. It is a truncation of the coherent family of fibres. A functor into $\Hot$ alone need not retain all that family's higher twisting, a distinction already illustrated in the commentary on Chapter~VI, manuscript pp.~7 and~74.

For a connected pointed map $f:(X',x'_0)\to(X,a)$, this transport has an explicit component formula. Set $G=\pi_1(X,a)$ and $H=\operatorname{im}\pi_1(f)$. Then
\[
                   \pi_0\operatorname{hofib}_a(f)\cong H\backslash G.
\]
To see the formula, join $x'_0$ to $x'$ by a path $\delta$. The pair $(x',\gamma)$ determines the loop $f\delta*\gamma$. A different $\delta$ multiplies its class on the left by an element of $H$. Conversely, a homotopy of pairs gives precisely this relation, by tracing the path of their first coordinates; any such relation is realized by that construction. A loop $c$ at $a$ therefore acts by right multiplication $Hg\mapsto Hg[c]$. For the inclusion $2\mathbb Z\subset\mathbb Z$, viewed as a map of one-object groupoids, there are two components and a generator of the base interchanges them. This answers the request about $\pi_0$ without choosing representatives in a categorical lifting category.

There is a functorial simplicial path object for every simplicial set $K$, not only for a Kan complex. Choose a functorial Kan replacement $K\to RK$ and put
\[
 P(K)=(K\times K)\times_{RK\times RK}(RK)^{\Delta^1}.
\]
Its endpoint map is a Kan fibration. The constant-path map $K\to P(K)$ is a weak equivalence: right properness first identifies $P(K)\to(RK)^{\Delta^1}$ as a weak equivalence, and evaluation identifies the latter with $RK$. Apply this with $K=NX$. Once $i:X\to\Ch_\infty(X)$ has been proved a weak equivalence, the span
\[
             N\Ch_\infty(X)\ \longleftarrow\ NX
                              \ \longrightarrow\ P(NX)
\]
is already a zigzag of weak equivalences over $NX\times NX$. This gives the first form of comparison in item~$5^o$. Identifying the \emph{strict} endpoint fibres with ordinary path spaces still requires the homotopy-fibre property of the categorical endpoint map; a weak equivalence of total spaces over a base does not by itself identify strict fibres.

For a general coefficient $A\in\mathbb D(e)$, the analogous homology transport is covariant; cohomology transport uses pullback along the inverse fibre equivalence. These constructions require the homotopy base-change isomorphisms appropriate to the stated $\mathbb D$-fibration.

The later theory of locally constant diagrams makes the distinction between a coherent family and its monodromy systematic \citep[Theorem~1.3 and the discussion preceding it]{Cisinski2009LocallyConstant}. The language of straightening provides a higher-categorical model of the family itself \citep[Theorem~2.2.1.2]{Lurie2009HTT}. These results explain the final form of the construction; the elementary fibre comparison remains the work that this chapter undertakes.
\end{dcomment}

\begin{dcomment}{exposition}{p144}{Why a category of lifts is the right test}
Grothendieck now replaces the geometric word ``smooth'' by a calculation that can be carried out on paths. Fix an arrow $w:b\to b'$ in a base $B$ and an object $e'$ above $b'$. For $p:E\to B$, the relevant category has objects $(e,\varphi)$ with $p(e)=b$ and $\varphi:e\to e'$ satisfying $p(\varphi)=w$. A morphism from $(e,\varphi)$ to $(e_1,\varphi_1)$ is a vertical map $h:e\to e_1$ with $\varphi_1h=\varphi$. Thus one must control both existence of lifts and the homotopy type of the space of choices.

This is exactly a comma category for the inclusion $E_b\to b\backslash E$. Requiring it to be contractible for every $(e',w)$ says that this inclusion is aspherical. The equivalence with smoothness is the criterion of \citet[Theorem~1.2]{Maltsiniotis2005Asphericite}. In the present case $B=X\times X$ and $w=(u,v)$. The annotation placing $\varphi$ inside $\Ch_\infty(X;a',b')$ cannot be literal: its source has endpoints $(a,b)$, so it is a morphism of the total path category over $(u,v)$.

The next construction is designed to make these comma categories manageable. At a finite interval whose two endpoint arrows point into the interior, a lift is obtained by changing the two endpoints and composing with $u$ and $v$. This is a Cartesian lift, hence a terminal object in the corresponding category of lifts. To deal with a path already lying in some stage, enlarge its support to a stage with the required endpoint parity. The filtered union then accommodates every finite diagram in the lifting category. This is why the two-index system introduced next is useful: it permits the same infinite path category to be approached by finite stages suited separately to incoming and outgoing lifts.

The proof of properness uses the opposite variance, as the duality formula on ms.~143 indicates. It does not follow from the mere existence of one lift, nor from a statement about the homotopy type of the total path category. The full calculation is carried out on ms.~150--153.
\end{dcomment}

\begin{dcomment}{exposition}{p147:La direction initiale est conforme}{Finite support and endpoint parity}
Grothendieck's two-sided sequence makes the eventual constancy of a path concrete. Its advantage is that support may be enlarged independently at the two ends. To avoid the shift in edge numbering between ms.~145 and ms.~146, write $e_i$ for the edge joining $a_i$ and $a_{i+1}$, directed forwards for even $i$ and backwards for odd $i$. Thus $e_i$ is the $u_{i+1}$ of the large diagram on ms.~145. A path supported on $[n,m]$ has $e_i$ equal to an identity outside $n\leq i<m$.

For $n<m$, the source-end sign is $(-1)^n$ and the target-end sign is $(-1)^{m-1}=(-1)^{m+1}$. The four cases are
\[
\begin{array}{c|c|c|c}
 n\bmod2&m\bmod2&\hbox{initial sign}&\hbox{terminal sign}\\ \hline
 0&0&+&-\\
 0&1&+&+\\
 1&0&-&-\\
 1&1&-&+
\end{array}
\]
For an odd length $m-n$, the number of full two-letter blocks is $(m-n-1)/2$; this is the correction already supplied by the editors in both displayed formulas. When $n=m$ there are no endpoint arrows and the subcategory is the constant-path copy of $X$; the sign table is not intended for that degenerate case.

The inclusion for $n'\leq n\leq m\leq m'$ extends a diagram by identity arrows. It is fully faithful, since a natural transformation is forced to be constant on each identity tail. Every eventually constant path and every natural transformation between two such paths belongs to a sufficiently large common support interval. Hence the filtered colimit of these full subcategories is the asserted category of two-sided paths, on morphisms as well as objects. The even-even intervals form a cofinal family of support intervals, and so do the odd-odd intervals. The first family has both endpoint arrows pointing into the interval and is adapted to Cartesian lifts; the second has both pointing out and is adapted to coCartesian lifts. Cofinality here is the ordinary directedness of support intervals, not an assertion that a fixed inclusion of categories is a weak equivalence.
\end{dcomment}

\begin{dcomment}{examples}{p148:de sorte que c'est le groupe semidirect de}{The dihedral action on a two-sided path}
Grothendieck introduces translations to compare the two parity choices without changing either endpoint. In the edge convention of the preceding block, translation and reversal are
\[
 (T^pa)_j=a_{j-p},\quad (T^pe)_j=e_{j-p},
 \qquad
 (Sa)_j=a_{-j},\quad (Se)_j=e_{-j-1}.
\]
These formulas apply equally to the components of a natural transformation, so they define functors, not merely permutations of the objects. Since paths are eventually constant, translation leaves both eventual values unchanged; reversal exchanges them. Odd translation reverses parity, whereas even translation preserves it. Reversal preserves each parity convention because reversing the direction of an edge also reverses the parity of its index.

The group relations are identities of functors. On vertices,
\[
 (S^2a)_j=a_j,
 \qquad
 (STSa)_j=a_{j+1}=(T^{-1}a)_j,
\]
and substitution in the edge formula gives the same equalities on edges and natural transformations. The acting group is therefore
\[
 D_\infty=\mathbb Z\rtimes\mathbb Z/2,
 \qquad (\epsilon,k):j\longmapsto\epsilon j+k,
 \quad \epsilon\in\{1,-1\}.
\]
The integer subgroup is the translation subgroup and reflection acts on it by inversion. This specifies the semidirect-product convention independently of the order of the groups in the French phrase.

For example, a path with one nonidentity edge $e_0:x\to y$ is carried by $T^2$ to the same edge at position $2$, with two additional identity positions on one side and two fewer on the other. Both have source $x$ and target $y$, but they are different objects of the strict path category. The formulas above prove an isomorphism of categories and the dihedral relations. They do not by themselves prove that $T^2$ is the identity after passage to $\Hot$. This is the specific coherence question taken up on the next page. It should be resolved by a homotopy or a comparison of the relevant diagrams, not by treating coordinate translation as literal equality.
\end{dcomment}


\begin{dcomment}{exposition}{p150}{Why both parities give the endpoint fibration}
Grothendieck wants the categorical path object to carry the same geometric information as an ordinary path fibration, before any fibrant replacement of $X$. The two ends of a finite zigzag already explain how this can happen. For a positive-length interval, when both endpoints are sources in the interval category, changing their values contravariantly only composes the two outer arrows; this gives a Grothendieck fibration over $X\times X$. When both endpoints are sinks, the analogous operation is covariant and gives an opfibration. The two parity choices, with positive length imposed, are cofinal among the finite supports of an eventually constant path. The zero-length case is the diagonal functor and must not be used in the fibration claim.

The passage to the filtered union deserves a proof, since the inclusion functors need not preserve Cartesian arrows. For a functor $q:E\to B$, an object $e\in E$, and an arrow $h:b\to q(e)$, put
\[
 L_q(e,h)=\{(d,\gamma):q(d)=b,\ \gamma:d\to e,\ q(\gamma)=h\},
\]
with morphisms over $1_b$ commuting with $\gamma$. Smoothness means that every such category is weakly contractible; properness has the dual outgoing-lift test. These are the fibre-to-comma criteria of \citet[Theorem~1.2 and its dual]{Maltsiniotis2005Asphericite}. Suppose $E=\varinjlim_i E_i$ over a filtered category, with all functors over $B$. Once $e$ is represented at one stage, its lift category is the filtered colimit of the corresponding lift categories at later stages. Every object, arrow, and equation defining a finite simplex occurs at a stage, and every equality holds at some later stage. Nerves preserve these filtered colimits. A filtered colimit of weakly contractible simplicial sets is weakly contractible: after $\operatorname{Ex}^{\infty}$, which also commutes with filtered colimits, the assertions about based homotopy groups involve only finitely many simplices. Consequently smoothness and properness pass separately to this filtered colimit. Preserving Cartesian arrows is unnecessary.

It follows that $\Ch_\infty(X)\to X\times X$ is smooth and proper. We can also justify directly the later marginal observation that these two conditions imply the required homotopy-fibre property. For $h:b\to b'$ in $B$, let $M_h$ be the category of arrows $d\to d'$ in $E$ over $h$, with its two projections to $E_b$ and $E_{b'}$. Properness and smoothness show that these projections are weak equivalences. For example, the comma category under an object $d_0\in E_b$ retracts, by precomposition, onto its outgoing-lift category; the dual argument uses postcomposition. Quillen's Theorem~A applies. The arrows defining $M_h$ then give a natural transformation between its two maps into $b\backslash E$, the second obtained by precomposition from $b'\backslash E$. Since smoothness identifies each strict fibre with its under-comma category, precomposition
\[
 b'\backslash E\longrightarrow b\backslash E
\]
is a weak equivalence. Quillen's Theorem~B identifies the strict fibres with homotopy fibres \citep[\S1, Theorem~B]{Quillen1973}. This proof also survives every base change. It establishes what the manuscript needs without the finite-test conservativity invoked on the following pages.
\end{dcomment}

\begin{dcomment}{examples}{p151}{A homotopy fibration without Cartesian lifts}
The warning in footnote~36 can be seen in the walking-arrow category $X=[1]$. It explains why a contractible category of choices is the appropriate conclusion, even when no best choice exists.

Consider the constant eventual path $c=1$ and the base arrow $(0\to1,1_1):(0,1)\to(1,1)$. An incoming lift is precisely an eventual alternating path $d$ from $0$ to $1$: its natural transformation to the constant path $1$ is unique. If the endpoint map were a Grothendieck fibration, this lifting category would have a terminal object. Its arrows are pointwise inequalities between the $0$--$1$ diagrams. For any candidate $d$, choose a direct edge far enough to the left that $d$ is still constantly $0$ there. Let $e$ change from $0$ to $1$ across that edge and remain $1$ afterwards. This is an allowed alternating path, because all its other edges are identities. At a vertex immediately after the chosen edge, $e$ is $1$ and $d$ is $0$, so there is no natural transformation $e\to d$. No terminal lift exists.

Dually, take the constant path $0$ and the outgoing base arrow $(1_0,0\to1)$. A cocartesian lift would give an initial object among paths from $0$ to $1$. For any candidate $d$, choose a single-step path whose direct transition is farther right than the point where $d$ has become constantly $1$. There is no arrow from $d$ to this later-step path. Thus no initial lift exists either. The incoming and outgoing lift categories are nevertheless weakly contractible by the filtered-stage proof. This is the exact distinction between the two homotopical lifting properties proved here and an ordinary fibration or opfibration.
\end{dcomment}

\begin{dcomment}{exposition}{p152}{The contraction with one endpoint fixed}
The next task is to make a path starting at $a$ disappear while keeping its starting point fixed. This is stronger than knowing merely that the interval has contractible nerve, because an arbitrary weak equivalence of interval categories need not induce a weak equivalence of strict functor categories. Here the interval has an explicit categorical deformation, and that is enough.

Write its vertices as $0,\ldots,r$, in their geometric order, with either orientation on each edge. For $0\leq k\leq r$, define $r_k(j)=\min(j,k)$. Each $r_k:I\to I$ is a functor. Successive $r_k,r_{k+1}$ are connected by a natural transformation in the direction of the edge between $k$ and $k+1$: on the tail $j\geq k+1$ use that edge, and elsewhere use identities. Naturality follows on the generating edges, hence on all arrows. Every $r_k$ fixes $0$, while $r_0$ is constant and $r_r=1_I$. Precomposition therefore joins the identity of $\operatorname{Fun}(I,X)_{a,?}$ to the constant functor at the constant path $a_I$, through a finite zigzag of natural transformations which is relative to the source. Its nerve is contractible.

Filtered colimits of these one-endpoint categories are weakly contractible by the finite simplicial argument just given. Thus both $\sigma$ and $\beta$ have weakly contractible strict fibres. Since they are smooth and proper, their pullbacks along arbitrary functors are weak equivalences. In the mapping-path construction, this universal statement is precisely what allows the category $X'$ to replace $X$ at the source endpoint. No supposition that $X$ itself is a Thomason fibrant category enters this argument.
\end{dcomment}

\begin{dcomment}{later}{p153}{Evrard's categorical cocylinder}
Grothendieck is seeking an explicit factorization of an arbitrary functor by adjoining paths. The same purpose already appears in \citet[\S I.1]{Evrard1975}, whose variable-length zigzags carry both natural transformations and maps which insert identity segments. The indexing conventions differ from those used here; an identification of the two categories would therefore require a comparison functor, not just the observation that their objects look alike.

A precise modern theorem is \citet[Theorem~2.8]{Shoikhet2016}: the categorical cocylinder gives a functorial factorization by a weak equivalence followed by a homotopy-fibred functor, and there is a dual construction. Its relation with the manuscript is the role of the second map: after applying the nerve, ordinary fibres compute homotopy fibres. Shoikhet's Proposition~1.14 formulates a stronger, base-change-stable version through simplicial sharp maps, under the prefibration hypotheses stated there. We use Quillen's comma-category theorem directly when those strict prefibration hypotheses have not been established. These distinctions matter for the total category $\Ch(X)$, whose type-changing arrows will obstruct properness below, even though its fixed-endpoint categories still have the correct homotopy types.
\end{dcomment}

\begin{dcomment}{exposition}{p156}{When the replacement is trivial}
The desired test for triviality is already formal once the mapping-path construction has been established. Put $r:C(f)\to X'$ for the projection which forgets the path. It is the pullback of the universally weakly equivalent source map $\sigma$, so $r$ is a weak equivalence. Its section $i_f$ is consequently a weak equivalence as well. The equality $\bar f i_f=f$ gives
\[
 \bar f\text{ is a weak equivalence}\quad\Longleftrightarrow\quad
 f\text{ is a weak equivalence}
\]
by two-out-of-three. If triviality is expressed instead by weakly contractible fibres, the homotopy-fibre theorem gives the same criterion. One can see both directions by replacing the realized map by a Serre fibration: a weak equivalence has contractible homotopy fibres, and a fibration with such fibres induces isomorphisms on every based homotopy group and a bijection on components.

Thus an independently developed long exact sequence for categorical loop objects is not needed for this particular conclusion. Nor should the preceding finite-test inference bear its proof. The claimed conservativity of the unbased restricted Yoneda functor is false in the generality used in the manuscript; the commentary on the finite-test obstruction in Chapter~V, manuscript p.~52, explains this. The two-out-of-three argument, together with the actual endpoint-fibration theorem, proves exactly the assertion required here.
\end{dcomment}

\begin{dcomment}{development}{p157}{A properness claim which the empty path disproves}
The return to paths of arbitrary type is motivated by their strictly associative concatenation. Their geometric usefulness is real, but the proposed proof gives them a stronger lifting property than they possess. We isolate the error before using the comparison diagram.

\textbf{The empty-path obstruction.}\ours\quad
For $X=[1]=(0\xrightarrow{w}1)$, neither
\[
 \Ch(X)\longrightarrow X\times X
 \qquad\text{nor}\qquad
 \mathcal C_{\mathbf N^o}\longrightarrow X\times X
\]
is proper, with the morphisms defined in this chapter.

\emph{Proof.} Let $z$ be the empty path at $0$, lying above $(0,0)$, and consider the base arrow $(1_0,w):(0,0)\to(0,1)$. Properness requires the category of its outgoing lifts with source $z$ to be weakly contractible. In $\Ch(X)$, a morphism from a path of type $e$ to a path of type $\tau$ lies over an endpoint-preserving type map $e\to\tau$. The single vertex of $e$ would have to map to both endpoints of $\tau$. Hence $\tau=e$. But an empty path has equal endpoints and cannot lie above $(0,1)$. The lift category is empty. For $\mathcal C_{\mathbf N^o}$ the same argument is even more immediate: an arrow from stage $0$ can only end at stage $0$, because its image in $\mathbf N^o$ has source $0$. Again its target must have equal endpoints. This proves both failures.\hfill$\square$

The opposite of a total Grothendieck construction reverses the arrows in its indexing category as well as the arrows inside its fibres. Reversing a path or replacing $X$ by $X^o$ does not perform that operation on $\Ty$ or on $\mathbf N^o$. This is the missing distinction in the deductions of properness on pp.~161 and 169. The same lifting test is the dual of \citet[Theorem~1.2]{Maltsiniotis2005Asphericite}.

This counterexample does not contradict the comparison of fixed-endpoint homotopy types which opens the section. That conclusion has a repair: smoothness of $\Ch(X)$, together with explicit cancellation of a path and its reverse, gives the homotopy-fibre theorem by Quillen~B. The intermediary $\mathcal C_{\mathbf N^o}$ can then be compared with $\Ch_\infty(X)$ using its filtered-colimit construction. We give the full argument below. We did not find this empty-path counterexample or this precise repair of the present pages in the primary sources checked, namely Evrard, Shoikhet, and Maltsiniotis; the mathematical ingredients of the repair are classical.
\end{dcomment}

\begin{dcomment}{exposition}{p160}{Completing the constant-path lemma}
The first comparison is a statement about total homotopy types, and here Grothendieck's argument works. Its useful feature is that it does not ask the chosen sequence of interval types to be homotopy cofinal in $\Ty$.

For every finite oriented interval $I$, evaluation at the initial vertex and the constant-diagram functor give
\[
 X\xrightarrow{c}\operatorname{Fun}(I,X)
 \xrightarrow{\operatorname{ev}_0}X,
 \qquad \operatorname{ev}_0c=1_X.
\]
The truncation functors $r_k(j)=\min(j,k)$ constructed above join $c\operatorname{ev}_0$ to the identity by a finite zigzag of natural transformations. This proves the lemma promised here and left unwritten on p.~170: $c$ is a categorical homotopy equivalence, not just a weak equivalence of nerves. The argument applies to every sign pattern and to the zero-length interval.

These constant-diagram functors commute strictly with precomposition by type maps. They therefore assemble to the displayed Cartesian functor
\[
 j:X\times\Ty\longrightarrow\Ch(X)
\]
between Grothendieck fibrations over $\Ty$, and each of its fibre functors is a weak equivalence. The standard fibrewise comparison theorem for Grothendieck constructions makes $j$ a weak equivalence. This is also the objectwise invariance behind Thomason's homotopy-colimit theorem; the fibred convention is recalled explicitly in \citet[\S1.5]{Shoikhet2016}. Since $\Ty$ has terminal object $e$, the inclusion $X\to X\times\Ty$ at $e$ is a weak equivalence. Their composite is the empty-path section $X\to\Ch(X)$.

Exactly the same proof over $\mathbf N^o$ proves $X\to\mathcal C_{\mathbf N^o}$ is a weak equivalence. For $\Ch_\infty(X)$ the constant-diagram functors assemble to a levelwise weak equivalence of filtered diagrams; invariance of weak equivalences under filtered colimits gives the corresponding assertion. Consequently the two comparison arrows out of $\mathcal C_{\mathbf N^o}$ are weak equivalences on total categories, by two-out-of-three. What remains is genuinely relative: a weak equivalence over $X\times X$ need not be a weak equivalence on its strict fibres until those fibres have been identified as homotopy fibres.
\end{dcomment}

\begin{dcomment}{exposition}{p164}{The lifting category is contracted relative to its endpoints}
The long diagram calculation has a simple purpose: to replace the category of all incoming lifts of $(u,v)$ by a slice with a terminal object. It is a proof of smoothness, and no properness assertion is needed to make it work.

Fix $u:a'\to a$, $v:b'\to b$, and a path $c:a\rightsquigarrow b$. Write $A=\Ch(X;a',b')$ and $\widetilde c=\mathrm{ch}(u)c\mathrm{ch}(v)^o$. Its slice $A/\widetilde c$ has terminal object $1_{\widetilde c}$. The functor $\varphi$ pads an incoming lift $w':c'\to c$ with a constant edge at each end. Its new map to $\widetilde c$ has the components $(1_{a'},u)$ on the left and $(v,1_{b'})$ on the right. A morphism $\alpha:c'\to c'_1$ is padded by identities, so the equality $w'_1\alpha=w'$ becomes $\widetilde w'_1\widetilde\alpha=\widetilde w'$. Both functoriality and the slice condition follow from the interchange law for concatenation of natural transformations.

To recover a lift from a slice object, there is an arrow $\rho_c:\widetilde c\to c$ over the collapse $\Delta_+\tau\Delta_-\to\tau$, with endpoint components $u,v$. Define $\psi(c'',w'')=(c'',\rho_cw'')$. The identity which makes the argument relative is
\[
 \rho_c\widetilde w'=w'\rho_{c'}.
\]
It holds on the central interval by the definition of $\widetilde w'$, and on the two added edges by the two commutative endpoint squares. The maps $\rho_{c'}$ are natural because collapsing constant collars commutes with every endpoint-preserving type map. They give a natural transformation $\psi\varphi\Rightarrow1$ on the lifting category.

On the other side, both $\varphi\psi$ and the identity map naturally to the constant endofunctor with value $1_{\widetilde c}$, because the slice has a terminal object. Thus their nerves are homotopic. The two functors are homotopy inverses, and the lifting category is contractible. This proves smoothness of $\Ch(X)\to X\times X$ completely. The wrong primes in the first padding diagram and the wrong final vertex in the diagram of $\rho_c$ are printed errors; the corrected endpoint descriptions above specify the maps without ambiguity.
\end{dcomment}

\begin{dcomment}{development}{p166}{Recovering the path-space theorem without properness}
The smoothness argument supplies enough control once we also know that transport along an arrow of $X\times X$ is invertible up to homotopy. Full path categories are particularly convenient for this second step, because their concatenation is a genuine functor.

\textbf{Transport and homotopy fibres for all path types.}\ours\quad
Let $E=\Ch(X)$ with endpoint functor $p:E\to B=X\times X$. For every $(a,b)$, the natural square with upper left corner $\Ch(X;a,b)$ and lower left corner $\{(a,b)\}$ is homotopy Cartesian. In particular,
\[
 B\Ch(X;a,b)\simeq
 \operatorname{hofib}_{(a,b)}(BX\xrightarrow{\Delta}BX\times BX),
\]
where $B(-)$ denotes the classifying space, so the right side is the ordinary path space from $a$ to $b$ up to weak equivalence.

\emph{Proof.} For $h=(u,v):(a',b')\to(a,b)$ define
\[
 S_h:\Ch(X;a,b)\longrightarrow\Ch(X;a',b'),
 \qquad c\longmapsto\mathrm{ch}(u)c\mathrm{ch}(v)^o.
\]
The opposite paths define a candidate inverse. Its composites insert the two backtracking paths associated to $u$ and $v$. For an arrow $u:x\to y$, cancellation is provided by the explicit diagrams
\[
 \text{empty}_x\ \longleftarrow\ x_{\Delta_+\Delta_-}
 \ \longrightarrow\ (x\xrightarrow{u}y\xleftarrow{u}x),
\]
\[
 (y\xleftarrow{u}x\xrightarrow{u}y)
 \ \longrightarrow\ y_{\Delta_-\Delta_+}
 \ \longrightarrow\ \text{empty}_y.
\]
The first rightward arrow has components $(1_x,u,1_x)$; the second display's first arrow has components $(1_y,u,1_y)$. The remaining arrows collapse constant intervals. They are morphisms with fixed endpoints. Whiskering with an arbitrary path makes these diagrams natural in that path. Cancelling each of the two inserted backtracks consequently gives homotopies from both composites to the identity. Hence $S_h$ is a homotopy equivalence.

Smoothness identifies the inclusion $i_z:E_z\to z\backslash E$ as a weak equivalence for every $z\in B$. The arrows $\rho_c:S_h(c)\to c$ over $h$ give a natural transformation
\[
 i_{(a',b')}S_h\ \Longrightarrow\ h^*i_{(a,b)}
\]
in $(a',b')\backslash E$. The three known weak equivalences in this square force $h^*:(a,b)\backslash E\to(a',b')\backslash E$ to be a weak equivalence. Quillen's Theorem~B now identifies $E_z$ with the homotopy fibre of $p$. Finally the constant-path section $X\to E$ is a weak equivalence over $X\times X$, so it identifies that homotopy fibre with the homotopy fibre of the diagonal. The latter is equivalent to the space of paths from $a$ to $b$: in the standard homotopy-fibre model, join the two paths from the chosen intermediate point, reversing the first.\hfill$\square$

The use of Quillen~B is classical \citep[\S1]{Quillen1973}; the explicit application above repairs the proof in these pages. We have checked the categorical cocylinder treatments of \citet{Evrard1975} and \citet{Shoikhet2016}, but did not locate this precise correction of Grothendieck's total path category there. No claim of priority is attached to the argument.
\end{dcomment}

\begin{dcomment}{development}{p169}{The sequential intermediary and the repaired comparison}
Grothendieck's intermediary remembers a finite stage as part of an object. That memory explains both its value for comparison and the limitation of the padding argument just attempted. Put $D=\Delta_+\Delta_-$, $t_n=D^{2n}$, and denote the specified symmetric-collapse transition by $\delta_{m,n}:t_m\to t_n$. An arrow of $\mathcal C_{\mathbf N^o}$ over $m\to n$ has exactly this underlying type map. Padding both ends of an arbitrary such arrow produces
\[
 1_D*\delta_{m,n}*1_D:t_{m+1}\longrightarrow t_{n+1},
\]
which is generally different from $\delta_{m+1,n+1}$. The former retains the outermost blocks and collapses blocks immediately inside them; the latter collapses the outermost blocks. Thus the formula on p.~168 need not even define the promised object of the slice of $\mathcal C_{\mathbf N^o}$. This is a restriction of the chosen indexing arrows, independent of the corrected exponents and primes.

\textbf{Smoothness of the sequential intermediary.}\ours\quad
Nevertheless $\mathcal C_{\mathbf N^o}\to X\times X$ is smooth.

\emph{Proof.} Fix a target path $c$ of stage $n$ and an incoming base arrow $h:(a',b')\to(a,b)$. Its lift category is fibred over the tail $\{m\geq n\}^o$. The fibre at $m$ is the incoming-lift category of $h$ into $\delta_{m,n}^*c$ for the fixed-type endpoint functor
\[
 \operatorname{Fun}(t_m,X)\longrightarrow X\times X.
\]
For $m>0$ this functor is a Grothendieck fibration: both endpoints of $t_m$ are sources, and a Cartesian lift is obtained by changing their objects and composing the first and last arrows with the specified endpoint maps. The corresponding lift category has that Cartesian lift as a terminal object. The only exception is $m=0$, possible when $n=0$. Its fibre is empty unless $a'=b'$ and the two components of $h$ coincide; when these equalities hold it is a singleton. If it is empty, omit stage $0$; otherwise retain it. In either case every fibre over the remaining tail is contractible, and that tail has a terminal object. The fibrewise comparison theorem for Grothendieck fibrations, applied to the projection to this tail, shows that the total lift category is weakly contractible. This is the smoothness criterion.\hfill$\square$

This argument proves only smoothness. Properness is excluded by the empty-path example above. Likewise the word ``asph\'erique'' at the end of the proof of $(1')$ on p.~166 is stronger than what was proved: for a positive-stage object $c$, the comma category of $X\to\mathcal C_{\mathbf N^o}$ over $c$ is empty, since there is no base arrow $0\to n$ for $n>0$. The section is a weak equivalence, as required.

\textbf{Comparison on every fixed-endpoint fibre.}\ours\quad
The canonical span
\[
 \Ch(X;a,b)\ \longleftarrow\
 \mathcal C_{\mathbf N^o}(a,b)\ \longrightarrow\
 \Ch_\infty(X;a,b)
\]
consists of weak equivalences.

\emph{Proof.} Write $E_1=\Ch(X)$, $E_0=\mathcal C_{\mathbf N^o}$, $E_2=\Ch_\infty(X)$, all over $B=X\times X$, and $\pi_i:E_0\to E_i$. The constant-path lemma proves both $\pi_i$ are weak equivalences on total categories. The filtered Grothendieck-to-colimit comparison, established in the commentary at manuscript pages~108--109, proves $\pi_2:E_{0,z}\to E_{2,z}$ is a weak equivalence on each strict fibre: apply that comparison to the sequence of fixed-endpoint categories itself.

By the endpoint theorem for $E_2$, its strict fibre is its homotopy fibre over $z$. The natural diagram
\[
\begin{tikzcd}[column sep=large]
 E_{0,z}\arrow[r,"\pi_2"]\arrow[d] & E_{2,z}\arrow[d,"\sim"]\\
 \operatorname{hofib}_z(BE_0\to BB)\arrow[r,"\sim"] &
 \operatorname{hofib}_z(BE_2\to BB)
\end{tikzcd}
\]
is understood after applying $B$ to the top row. Its lower arrow is a weak equivalence because the maps on total spaces are weak equivalences over the same base. The left vertical map is therefore a weak equivalence. The preceding block proves the same homotopy-fibre description for $E_1$. Repeating the square with $\pi_1$ shows that $E_{0,z}\to E_{1,z}$ is a weak equivalence.\hfill$\square$

The conclusion of the section is thus recovered, although its claimed properness and its particular sequential padding construction cannot be retained. The repair uses based homotopy fibres, never a conservativity assertion about unbased finite tests. The literature check is the one stated in the preceding development block; these are proofs supplied here for the specified models.
\end{dcomment}

\begin{dcomment}{development}{p169}{Even translation acts trivially after localization}
The comparison now also settles the question about $T^2$ left on pp.~149--150. Its resolution uses a visible change of finite support, rather than equality after restriction to finite test objects.

\textbf{Translation invariance.}\ours\quad
The even translation of an eventually constant alternating path induces the identity in $\Hot$, both on $\Ch_\infty(X)$ and on each fixed-endpoint category.

\emph{Proof.} Keep $D=\Delta_+\Delta_-$ and $t_n=D^{2n}$. Define $R$ on the sequential Grothendieck category by sending a path of stage $n$ to stage $n+1$, adding two constant $D$-blocks at its right end. The asymmetric collapse $q_n:t_{n+1}\to t_n$ defining this operation commutes with the symmetric transitions:
\[
 q_n\delta_{m+1,n+1}=\delta_{m,n}q_m\qquad(m\geq n).
\]
Both sides collapse the same number of left and right blocks. Hence $R$ is a functor. On the eventual bi-infinite realization, with the centred indexing of these pages this operation is $U=T^{-2}$, and $\pi_2R=U\pi_2$.

The maps $q_n$ are permitted in $\Ty$. They give a natural transformation $\pi_1R\Rightarrow\pi_1$ into the full path category, whose component is the Cartesian arrow from the padded path to the original path. By the comparison just proved, $k=\pi_1\pi_2^{-1}$ is an isomorphism in $\Hot$. Therefore
\[
 kU\pi_2=\pi_1R=\pi_1=k\pi_2,
\]
and cancellation gives $U=1$. The same construction fixes the endpoints, so the argument applies to every fibre. Taking inverses gives $T^2=1$ as well.\hfill$\square$

We did not find this finite-support proof in the checked literature. It is a consequence of the repaired comparison, with no additional general theorem about the restricted Yoneda functor.
\end{dcomment}

\begin{dcomment}{exposition}{p170}{The formal diagonal factorization, with its hypothesis retained}
The closing argument extracts the reusable part of the path construction. Suppose a category $C_X$ is given with a factorization
\[
 X\xrightarrow{i_X}C_X\xrightarrow{p_X}X\times X,
 \qquad p_Xi_X=\Delta_X,
\]
where $i_X$ is a weak equivalence and $p_X$ is smooth and proper. Put $\sigma_X=\operatorname{pr}_1p_X$ and $\beta_X=\operatorname{pr}_2p_X$. Both are smooth and proper, because these classes are closed under composition and base change; a projection is a base change of a map to the terminal category. The identity $\sigma_Xi_X=1_X$ implies that $\sigma_X$ is a weak equivalence. Its strict fibres are homotopy fibres by the endpoint argument, hence weakly contractible. The lift-category criterion and Quillen~A then show that every base change of $\sigma_X$ is a weak equivalence.

For $f:X'\to X$, form the ordinary pullback $C(f)=X'\times_{X,\sigma_X}C_X$. Its projection $r:C(f)\to X'$ is a weak equivalence. The section $i_f$ induced by $i_X$ satisfies $ri_f=1_{X'}$, so $i_f$ is a weak equivalence. The map to the other endpoint factors as
\[
 C(f)\longrightarrow X'\times X\longrightarrow X.
\]
The first arrow is the base change of $p_X$, and the second is a projection. Their composite $\bar f$ is therefore smooth and proper. The equation $\bar f i_f=f$ proves the requested factorization. This is exactly the formal argument of the manuscript, with the universal weak-equivalence step justified.

It applies to $C_X=\Ch_\infty(X)$, whose two lifting properties were proved using the two cofinal parity systems. It cannot be applied with the asserted hypothesis to $C_X=\Ch(X)$ or $\mathcal C_{\mathbf N^o}$: those endpoint maps fail properness. Their weak equivalence over $X\times X$ does not transfer a strict comma-category lifting condition. The formal result and the comparison theorem therefore coexist without identifying their hypotheses.
\end{dcomment}

\begin{dcomment}{examples}{p171}{A mapping-path fibre for a homomorphism of groups}
The factorization records more than the kernel of a functor on components. For a homomorphism $\phi:H\to G$, write $\mathbf B H$ and $\mathbf B G$ for their one-object groupoids and consider the functor $\mathbf B\phi$ between them. The homotopy fibre has a particularly small categorical model, the action groupoid $F_\phi$ with objects $g\in G$ and arrows
\[
 h:g\longrightarrow g'
 \quad\Longleftrightarrow\quad
 g'\phi(h)=g.
\]
If $h:g\to g'$ and $k:g'\to g''$, their composite is $kh$, since $g''\phi(kh)=g$. All arrows are invertible. This is the over-comma category $(\mathbf B\phi)/*$, where $*$ is the unique object of $\mathbf B G$. Every arrow in the base groupoid acts on it by an isomorphism, so Quillen~B identifies $BF_\phi$ with the homotopy fibre.

Two objects lie in the same component precisely when they are in the same coset $g\operatorname{im}\phi$. An automorphism of any $g$ is an $h$ with $\phi(h)=1$, hence the isotropy group is $\ker\phi$. Choosing one representative in each coset gives an equivalence of groupoids
\[
 F_\phi\simeq\coprod_{G/\operatorname{im}\phi}\mathbf B(\ker\phi).
\]
Thus $\pi_0$ detects the failure of surjectivity and the fundamental group of each component detects the failure of injectivity. For multiplication by a nonzero integer $m$ on $\mathbf Z$, the fibre has $|m|$ contractible components. For the map $\mathbf Z\to1$, it is a circle up to homotopy. The corresponding fibre of Grothendieck's mapping-path factorization has exactly these homotopy types. In particular it is weakly contractible precisely when $\phi$ is an isomorphism, which illustrates the triviality criterion at p.~156 without computing any individual categorical path.
\end{dcomment}


\begin{dcomment}{development}{p169}{Why the comparison preserves composition}
Grothendieck's original reason for retaining all finite path types was their strictly associative concatenation. The comparison of endpoint fibres should therefore identify that operation with ordinary path concatenation, not merely identify the underlying homotopy types. The constant-path sections determine this compatibility once composable pairs are treated as carefully as individual paths.

\textbf{Compatibility with path concatenation.}\ours\quad
Put $E=\Ch(X)$ and $K=BX$. The equivalences
\[
 u_{a,b}:B\Ch(X;a,b)\longrightarrow P_{a,b}K
 \quad\hbox{in }\Hot,
\]
obtained from the constant-path diagonal factorization, can be chosen canonically in the localized slice category and satisfy
\[
 u_{a,c}\,B(\mathord{*})
   =\mathord{*}_{\mathrm{top}}(u_{a,b}\times u_{b,c})
 \quad\hbox{in }\Hot.
\]
Here $P_{a,b}K$ is the ordinary space of paths from $a$ to $b$, and the products have the indicated endpoint fibres as their domain. In particular the equivalence on loops respects the usual multiplication in the homotopy category.

\emph{Proof.} Let $D=E\times_{X}E$, using the target of the first path and the source of the second. Its projection to $\Ty\times\Ty$ is the fibred Grothendieck construction with fibre
$\operatorname{Fun}(\tau*\upsilon,X)$ at $(\tau,\upsilon)$: matching the two endpoint values is exactly the condition for gluing the diagrams on the two intervals. The constant-diagram functor from $X$ to this fibre is a categorical homotopy equivalence, by the truncation homotopies proved at manuscript page~160. These functors commute with the type maps. The same fibrewise comparison used there, and the terminal object $(e,e)$ of $\Ty\times\Ty$, show that the constant-pair section $j:X\to D$ is a weak equivalence.

We also need the relative assertion for $D$. Its three-endpoint functor $r:D\to X^3$ factors as
\[
 D\longrightarrow E\times X\longrightarrow X^3,
 \qquad
 (c,d)\longmapsto(c,\beta d)
       \longmapsto(\sigma c,\beta c,\beta d).
\]
The first functor is the base change of $E\to X^2$ along $(c,z)\mapsto(\beta c,z)$; the second is its product with $1_X$. Thus $r$ is smooth, by stability under base change and composition \citep[Proposition~4.7]{Maltsiniotis2005Asphericite}. For a base arrow $(u,v,w):(a',b',c')\to(a,b,c)$ the functor on strict fibres is
\[
 S_{(u,v)}\times S_{(v,w)}.
\]
Both factors are homotopy equivalences by the explicit backtrack cancellations at page~166. The two collapse arrows $\rho$ agree on their common middle endpoint, with component $v$, so together give a natural arrow over $(u,v,w)$. Smoothness identifies each strict fibre with its under-comma category. The natural arrow then shows that every transition between those comma categories is a weak equivalence, exactly as in the proof for $E$. Quillen's Theorem~B applies. Consequently
\[
 B\bigl(\Ch(X;a,b)\times\Ch(X;b,c)\bigr)
     \longrightarrow\operatorname{hofib}_{(a,b,c)}(BD\to K^3)
\]
is a weak equivalence.

Now use ordinary spaces, with compactly generated mapping spaces. Put $P=K^{[0,1]}$ and $Q=P\times_K P$, and let $s:K\to P$ and $t:K\to Q$ be their constant-path sections. Both are weak equivalences: $P$ and $Q$ are the spaces of maps from the intervals $[0,1]$ and $[0,2]$, respectively, and evaluation at the initial point is a deformation retraction onto constants. Their endpoint maps, to $K^2$ and $K^3$, are fibrations. Write $i:K\to BE$ for the constant-path section and still write $j:K\to BD$ for its two-path version. In the homotopy categories of spaces over these endpoint bases, define
\[
 u=s\,i^{-1}:BE\longrightarrow P,
 \qquad
 v=t\,j^{-1}:BD\longrightarrow Q.
\]
These are the unique comparison isomorphisms whose composites with the constant sections are $s$ and $t$.

All maps in the remaining argument can be regarded over $K^3$. For concatenation the target is the pullback of $P\to K^2$ along $(x,y,z)\mapsto(x,z)$; for either projection it is the corresponding pullback along $(x,y,z)\mapsto(x,y)$ or $(y,z)$. These pullbacks are just products with the unused copy of $K$, so they preserve weak equivalences and the displayed localized comparisons. Both composites comparing categorical concatenation with topological concatenation send $j$ to the same constant path. Since $j$ is invertible in this slice homotopy category, cancellation gives
\[
                u\,B(\mathord{*})=\mathord{*}_{\mathrm{top}}\,v.
\]
The identical argument with either path projection shows that $v$ commutes with both projections under $u$.

Take the derived fibre over $(a,b,c)$. The relative assertion for $D$, the endpoint theorem for $E$, and the ordinary path fibrations identify these derived fibres with the stated strict path spaces and their products. The two projection identities imply that the induced map of the product fibres is $u_{a,b}\times u_{b,c}$, since a map to a product in $\Hot$ is determined by its two projections. The concatenation identity is therefore the required formula. Constants give the corresponding unit identity.\hfill$\square$

This argument supplies the compatibility omitted by a bare comparison of homotopy types. It uses Quillen's relative fibre theorem and localization at weak equivalences, not a conservativity or faithfulness assertion about finite tests. We checked the path constructions in \citet[Section~I.1]{Evrard1975} and the diagonal and cocylinder comparisons in \citet[Sections~1.5 and~2.1, Theorem~2.8]{Shoikhet2016}; we did not find this particular comparison through the category of composable pairs there. The ingredients are classical, and no priority is claimed.
\end{dcomment}
